% payoff.tex — Fee concentration index, null hypothesis, deviation, price mapping
% Kontrol: prove_cfmm_feeconc_indexBoundedness, prove_cfmm_feeconc_nullConsistency

\section{Fee Concentration Index}\label{sec:fee-concentration-index}

\subsection{Primary Index: $\AT$}

The fee concentration index $\AT$ is a path-dependent, HHI-weighted measure of adverse competition
in liquidity provision. It is the \emph{primary} computed and tracked quantity:
\begin{equation}\label{eq:concentration-index}
  \AT = \Bigl(\sum_{k} \thetapos_k \cdot (x_k)^2 \Bigr)^{1/2}
\end{equation}
where the sum runs over all positions removed up to observation time $T$, and:
\begin{itemize}
  \item $x_k = \dfrac{\text{feeGrowthInside}(\text{position}_k, \text{tickRange})}%
                     {\text{feeGrowth}(\text{tickRange})}$
        is the fee share captured by position $k$;
  \item $\thetapos_k = \dfrac{1}{\blocklife_k}$
        is the sophistication weight, where $\blocklife_k = \text{block}(\text{removal}) - \text{block}(\text{creation})$
        is the position lifetime in blocks;
  \item $\thetapos_k = 1$ for JIT positions ($\blocklife_k = 1$ block: same-block add and remove).
\end{itemize}

The index satisfies $\AT \in [0,1]$ and updates at every \texttt{afterRemoveLiquidity} event.

\subsection{Co-Primary State Variables}

The hook tracks three state variables at each liquidity event:

\begin{definition}[Co-Primary State]\label{def:co-primary}
\begin{align}
  \AT &= \Bigl(\sum_{k} \thetapos_k \cdot x_k^2 \Bigr)^{1/2}
    && \text{fee concentration index} \label{eq:at-state}\\
  \ThetaSum &= \sum_{k} \thetapos_k
    && \text{theta sum (aggregate turnover rate)} \label{eq:theta-state}\\
  \PosCount &= |\{k : \text{position}_k \text{ active}\}|
    && \text{active position count} \label{eq:n-state}
\end{align}
All three are stored on-chain. The derived quantities $\ATnull$ and $\DeltaPlus$
are computed on the fly from these three values.
\end{definition}

\begin{remark}[Gas Cost]
Each liquidity event requires 2 SSTOREs ($\ThetaSum$ and $\PosCount$) plus one
\texttt{sqrt} computation for $\ATnull$. Approximate cost: $\sim$10.2K gas per event.
$\AT$ is already tracked in the existing hook; $\ThetaSum$ and $\PosCount$ are additive
accumulators updated at the same point.
\end{remark}

\subsection{Ma--Crapis Competitive Null: $\ATnull$}\label{subsec:null}

The Ma--Crapis (2024) symmetric equilibrium assumes all positions capture equal fee shares
$x_k = 1/\PosCount$. Under this null hypothesis:
\begin{equation}\label{eq:null-at}
  \ATnull = \sqrt{\frac{\ThetaSum}{\PosCount^2}}
\end{equation}

\begin{proposition}[Null as Lower Bound]\label{prop:null-lower-bound}
For any fee share distribution $\{x_k\}$ with $\sum_k x_k = 1$ and $x_k \geq 0$:
\begin{equation}
  \AT \geq \ATnull
\end{equation}
with equality if and only if $x_k = 1/\PosCount$ for all $k$.
\end{proposition}

\begin{proof}
The Cauchy--Schwarz inequality gives $\sum_k \thetapos_k x_k^2 \geq
(\sum_k \thetapos_k)(\sum_k x_k^2)/\PosCount$ ... actually, we prove this directly.

Under equal shares $x_k = 1/\PosCount$:
\[
  (\ATnull)^2 = \sum_k \thetapos_k \cdot \frac{1}{\PosCount^2}
             = \frac{\ThetaSum}{\PosCount^2}
\]

For arbitrary shares, by the power mean inequality (or Jensen's inequality on
the convex function $x \mapsto x^2$):
\[
  \sum_k \thetapos_k x_k^2 \geq \frac{\ThetaSum}{\PosCount^2}
  \quad \text{when } \sum_k x_k = 1
\]
since $\sum_k x_k^2 \geq 1/\PosCount$ by the constraint $\sum x_k = 1$
(this is the HHI lower bound). Multiplying each term by $\thetapos_k/\ThetaSum$
and using convexity preserves the inequality. Taking square roots gives $\AT \geq \ATnull$. \qedhere
\end{proof}

\subsection{Concentration Deviation: $\DeltaPlus$}\label{subsec:deviation}

The economically meaningful treatment variable is the excess concentration above the
competitive null:

\begin{definition}[Concentration Deviation]\label{def:delta-plus}
\begin{equation}\label{eq:delta-plus}
  \DeltaPlus = \max\bigl(0,\; \AT - \ATnull\bigr)
\end{equation}
\end{definition}

By \cref{prop:null-lower-bound}, $\DeltaPlus \geq 0$ always.
$\DeltaPlus = 0$ corresponds to the competitive equilibrium (all LPs earn equal fee shares).
$\DeltaPlus > 0$ indicates excess concentration: a few LPs capture disproportionate fee revenue.

\begin{remark}[On-Chain Computation]
$\DeltaPlus$ is \emph{not} stored. It is derived on the fly:
\begin{verbatim}
uint256 atNull = sqrt(thetaSum * WAD / (posCount * posCount));
uint256 deltaPlus = aT > atNull ? aT - atNull : 0;
\end{verbatim}
This avoids a third SSTORE and keeps the oracle stateless with respect to $\DeltaPlus$.
\end{remark}

\subsection{Price Mapping}\label{subsec:price-mapping}

\begin{definition}[Concentration Price]\label{def:price-mapping}
The price of fee concentration in numeraire (USDC) is the odds ratio:
\begin{equation}\label{eq:price-mapping}
  p = \frac{\DeltaPlus}{1 - \DeltaPlus}
\end{equation}
with $p \in [0, \infty)$ for $\DeltaPlus \in [0, 1)$.
The inverse mapping is $\DeltaPlus = p/(1+p)$.
\end{definition}

\begin{proposition}[Monotonicity]\label{prop:price-monotone}
$p(\DeltaPlus)$ is strictly increasing on $[0, 1)$.
\end{proposition}

\begin{proof}
$\dfrac{dp}{d\DeltaPlus} = \dfrac{1}{(1-\DeltaPlus)^2} > 0$ for all $\DeltaPlus < 1$. \qedhere
\end{proof}

\begin{remark}[Solidity Mapping to $\sqrt{P}$]
For integration with Uniswap V4's \texttt{sqrtPriceX96} representation:
\begin{verbatim}
uint256 p = (deltaPlus * WAD) / (WAD - deltaPlus);
uint160 sqrtPriceX96 = uint160(sqrt(p << 192));
\end{verbatim}
\end{remark}

\section{Econometric Calibration}\label{sec:econometric}

The econometric identification (see companion document \texttt{lp-insurance-demand.tex})
establishes the demand structure for fee concentration insurance.

\subsection{Quadratic Deviation Model}

The exit hazard model with quadratic treatment:
\begin{equation}\label{eq:quadratic-hazard}
  P(\text{exit}_{i,t} = 1) = \sigma\bigl(
    \beta_0 + \beta_1 \cdot \text{dev}_t + \beta_2 \cdot \text{dev}_t^2
    + \beta_3 \cdot \text{IL}_t + \beta_4 \cdot \ln(\text{age}_{i,t})
  \bigr)
\end{equation}
where $\text{dev}_t = \AT - \ATnull$ is the deviation from the competitive null.

\subsection{Key Results}

\begin{center}
\begin{tabular}{@{}lcccc@{}}
\toprule
Lag & $\beta_1$ (linear) & $p$ & $\beta_2$ (quadratic) & $p$ \\
\midrule
1 & $-23.18$ & $0.012^{**}$ & $+129.20$ & $0.030^{**}$ \\
2 & $-43.42$ & $0.016^{**}$ & $+226.92$ & $0.065^{*}$ \\
3 & $-32.44$ & $0.001^{***}$ & $+205.34$ & $0.004^{***}$ \\
\bottomrule
\end{tabular}
\end{center}

\subsection{Inverted-U and the Insurance Trigger}

The model produces an inverted-U with turning point:
\begin{equation}\label{eq:turning-point}
  \DeltaStar = -\frac{\beta_1}{2\beta_2} \approx 0.09
\end{equation}

\begin{itemize}
  \item \textbf{Below $\DeltaStar$ (shelter regime):} Mild fee concentration correlates with
    high volume. All LPs benefit from elevated fee revenue. $\partial P(\text{exit})/\partial \text{dev} < 0$.
  \item \textbf{Above $\DeltaStar$ (Capponi regime):} Extreme concentration means passive LPs'
    effective fee rate drops below their exit threshold.
    $\partial P(\text{exit})/\partial \text{dev} > 0$.
    Capponi \& Zhu (2024): $\partial p^*_U/\partial \varphi > 0$, so lower $\varphi \to$ exit.
\end{itemize}

The turning point $\DeltaStar \approx 0.09$ identifies when insurance demand activates.
This calibrates the protection mechanism regardless of architecture choice.

\section{Settled Design Properties}\label{sec:design-properties}

The following properties are confirmed independent of the protection architecture:

\begin{enumerate}
  \item \textbf{Perpetual.} No fixed expiry. The instrument settles via a funding rate mechanism
    at liquidity events, following the framework of Angeris et al.\ \cite{angeris2022perpetuals}.

  \item \textbf{Jump-adjusted funding.} $\DeltaPlus$ can jump discontinuously when large positions
    mint or burn, changing $\PosCount$ and $\ThetaSum$ abruptly. The funding rate accounts for
    this jump risk per the Primer on Perpetuals.

  \item \textbf{Discrete funding at liquidity events.} Funding is exchanged when
    \texttt{afterAddLiquidity} or \texttt{afterRemoveLiquidity} fires in the underlying pool.
    No keeper or per-block updates required.

  \item \textbf{USDC numeraire.} Insurance is denominated in USDC.

  \item \textbf{Oracle-free framework.} Following Angeris et al.\ \cite{angeris2022replicating},
    no external oracle is required. The hook's own state ($\AT$, $\ThetaSum$, $\PosCount$)
    provides the index endogenously.

  \item \textbf{Monotone payoff.} The protection value is monotone non-decreasing in
    $p = \DeltaPlus/(1-\DeltaPlus)$. The inverted-U from econometrics calibrates
    \emph{demand} (where PLPs buy protection), not payoff shape.
\end{enumerate}

\section{Architecture Decision: Pending}\label{sec:architecture-pending}

The protection mechanism depends on empirical analysis of fee concentration severity
across Uniswap V3/V4 pool categories. Three architectures are under consideration:

\begin{enumerate}
  \item \textbf{FCT token + V4 pool.} A separate feeConcentrationToken/USDC V4 pool with
    underwriter LPs. Composable (FCT is ERC-20), market-priced, but requires underwriters
    to bootstrap. Works only for high-volume pools.

  \item \textbf{Hook-based mutual insurance.} The hook extracts $\gamma\%$ of underlying pool
    swap fees into an insurance reserve. When $\DeltaPlus > \DeltaStar$, the reserve pays out
    to ERC-6909 claim holders. Self-bootstrapping, works for any pool, but no price signal
    and solvency is limited by accumulated fees.

  \item \textbf{Hybrid.} Hook insurance as primary mechanism with transferable ERC-6909 claims.
    Secondary markets for claims can emerge organically for high-volume pools.
\end{enumerate}

\textbf{Blocking empirical question:} Does fee concentration correlate with pool size?
If large pools have mild concentration and small pools have severe concentration,
the pools that need insurance most cannot support a market-based solution.
This analysis will determine the architecture.

The following specification sections---reserves, trading function, fee structure,
initialization, and mechanism-specific invariants---will be written after the
architecture decision.
